% !TEX TS-program = pdflatex
% !TEX encoding = UTF-8 Unicode

% This file is a template using the "beamer" package to create slides for a talk or presentation
% - Giving a talk on some subject.
% - The talk is between 15min and 45min long.
% - Style is ornate.

% MODIFIED by Jonathan Kew, 2008-07-06
% The header comments and encoding in this file were modified for inclusion with TeXworks.
% The content is otherwise unchanged from the original distributed with the beamer package.

\documentclass{beamer}
\setbeamercovered{invisible}



% Copyright 2004 by Till Tantau <tantau@users.sourceforge.net>.
%
% In principle, this file can be redistributed and/or modified under
% the terms of the GNU Public License, version 2.
%
% However, this file is supposed to be a template to be modified
% for your own needs. For this reason, if you use this file as a
% template and not specifically distribute it as part of a another
% package/program, I grant the extra permission to freely copy and
% modify this file as you see fit and even to delete this copyright
% notice. 


\mode<presentation>
{
  \usetheme{Warsaw}
  % or ...

  % or whatever (possibly just delete it)
}


\usepackage[english]{babel}
% or whatever

\usepackage[utf8]{inputenc}
% or whatever

\usepackage{times}
\usepackage[T1]{fontenc}
% Or whatever. Note that the encoding and the font should match. If T1
% does not look nice, try deleting the line with the fontenc.

%%% MATH RELATED

\input{macros.tex}

\title{MATH 350-2 Advanced Calculus} 
\subtitle
{} % (optional)

\author[W.R. Casper] % (optional, use only with lots of authors)
{W.R. Casper}
% - Use the \inst{?} command only if the authors have different
%   affiliation.

\institute[California State University Fullerton] % (optional, but mostly needed)
{
  Department of Mathematics\\
  California State University Fullerton}
% - Use the \inst command only if there are several affiliations.
% - Keep it simple, no one is interested in your street address.

\subject{Talks}
% This is only inserted into the PDF information catalog. Can be left
% out. 



% If you have a file called "university-logo-filename.xxx", where xxx
% is a graphic format that can be processed by latex or pdflatex,
% resp., then you can add a logo as follows:

% \pgfdeclareimage[height=0.5cm]{university-logo}{university-logo-filename}
% \logo{\pgfuseimage{university-logo}}



% Delete this, if you do not want the table of contents to pop up at
% the beginning of each subsection:
\AtBeginSubsection[]
{
  \begin{frame}<beamer>{Outline}
    \tableofcontents[currentsection,currentsubsection]
  \end{frame}
}


% If you wish to uncover everything in a step-wise fashion, uncomment
% the following command: 

%\beamerdefaultoverlayspecification{<+->}


\begin{document}

\begin{frame}
  \titlepage
\end{frame}

\begin{frame}{Outline}
  \tableofcontents
  % You might wish to add the option [pausesections]
\end{frame}

% Since this a solution template for a generic talk, very little can
% be said about how it should be structured. However, the talk length
% of between 15min and 45min and the theme suggest that you stick to
% the following rules:  

% - Exactly two or three sections (other than the summary).
% - At *most* three subsections per section.
% - Talk about 30s to 2min per frame. So there should be between about
%   15 and 30 frames, all told.

\section{Real Analysis Lecture 22}

\subsection{Riemann-Stieltjes integrals}

\begin{frame}{Stieltjes sums}
Let $f$ and $\alpha$ be real-valued functions on $[a,b]$, and let $P=\{x_0,x_1,\dots,x_n\}$ be a partition of $[a,b]$ and define\\
{\small
$$m_k(f) = \inf\{f(x): x_{k-1}\leq x\leq x_k\},\ \ M_k(f) = \sup\{f(x): x_{k-1}\leq x\leq x_k\}.$$
}
\pause
The \textbf{upper Stieltjes sum} of $f$ with respect to $\alpha$ on $[a,b]$ is:
$$\overline{S}(P,f,\alpha) = \sum_{k=1}^n M_k(f)\Delta\alpha_k(P),$$
\pause
and the \textbf{lower Stieltjes sum} of $f$ with respect to $\alpha$ on $[a,b]$ is:
$$\underline{S}(P,f,\alpha) = \sum_{k=1}^n m_k(f)\Delta\alpha_k(P)$$
\end{frame}

\begin{frame}{Stieltjes sums}
These sums refine predictably!
\begin{lem}
Let $f$ and $\alpha$ be real-valued functions on $[a,b]$, with $\alpha$ monotone increasing.
\pause
If $P$ is a partition of $[a,b]$ and $P'$ is a refinement of $P$, then
$$
\underline{S}(P,f,\alpha) \leq
\underline{S}(P',f,\alpha) \leq
\overline{S}(P',f,\alpha) \leq
\overline{S}(P,f,\alpha)
$$
\end{lem}
\end{frame}

\begin{frame}{Stieltjes sums}
\begin{proof}
The inequality $\underline{S}(P',f,\alpha) \leq \overline{S}(P',f,\alpha)$ is clear.\\
\pause
We will prove $\underline{S}(P,f,\alpha) \leq \underline{S}(P',f,\alpha)$.\\
\pause
Let $P=\{x_0,\dots, x_n\}$ and $P'=\{x_0',\dots, x_N'\}$.\\
\pause
For each $k$, let $\iota(k)$ with $x_{\iota(k)-1} \leq x_k'\leq x_{\iota(k)}$.\\
\pause
{
\tiny
$$m_{\iota(k)}(f)=\inf\{f(x): x_{\iota(k)-1}\leq x\leq x_{\iota(k)}\}\leq \inf\{f(x): x_{k-1}'\leq x\leq x_k'\}=m_k'(f)$$
}
\pause
Therefore 
{\tiny
\begin{align*}
\underline{S}(P',f,\alpha)
 &=    \sum_{k=1}^N m_k'(f)\Delta\alpha_k(P')\\
 &\geq \sum_{k=1}^N m_{\iota(k)}(f)\Delta\alpha_k(P')\\
 & = \sum_{k=1}^n m_{k}(f)\Delta\alpha_k(P) = \underline{S}(P,f,\alpha)
\end{align*}
}
\end{proof}
\end{frame}

\begin{frame}{Stieltjes integrals}
Let $f$ and $\alpha$ be real-valued functions on $[a,b]$, with $\alpha$ monotone increasing.
\pause
The \textbf{upper Stieltjes integral} of $f$ with respect to $\alpha$ on $[a,b]$ is:
$$\overline{\int} fd\alpha = \inf\{\overline{S}(P,f,\alpha): P\ \text{partition of}\ [a,b]\}$$
\pause
The \textbf{lower Stieltjes integral} of $f$ with respect to $\alpha$ on $[a,b]$ is:
$$\underline{\int} fd\alpha = \sup\{\underline{S}(P,f,\alpha): P\ \text{partition of}\ [a,b]\}$$
\end{frame}

\begin{frame}{Riemann's criterion}
\begin{defn}
Let $f$ and $\alpha$ be real-valued functions on $[a,b]$.\\
Then $f$ \textbf{satisfies Riemann's criterion} with respect to $\alpha$ on $[a,b]$ if for all $\epsilon > 0$ there exists a partition $P_\epsilon$ such that for all refinements $P$ of $P_\epsilon$
$$0\leq \overline S(P,f,\alpha)-\underline S(P,f,\alpha) < \epsilon.$$
\end{defn}
\pause
\begin{thm}
Let $f$ and $\alpha$ be real-valued functions on $[a,b]$, with $\alpha$ monotone increasing.  Then the following are equivalent:\\
\pause
(a) $f$ is Riemann integrable with respect to $\alpha$ on $[a,b]$\\
\pause
(b) $f$ satisfies Riemann's criterion with respect to $\alpha$ on $[a,b]$\\
\pause
(c) $\underline \int f d\alpha = \overline \int f d\alpha$
\end{thm}
\end{frame}

\begin{frame}{Challenge!}
Suppose that $f(x) = x$, $\alpha(x) = x$, and $[a,b] = [0,1]$.\\
\begin{prob}
Calculate the value of
$$\overline S(P,f,\alpha)\ \ \text{for}\ \ P = \{0/n,1/n,\dots,n/n\}.$$
\end{prob}
\end{frame}

\begin{frame}{Challenge!}
Suppose that $f(x) = x$, $\alpha(x) = x$, and $[a,b] = [0,1]$.\\
\begin{prob}
Calculate the value of
$$\underline S(P,f,\alpha)\ \ \text{for}\ \ P = \{0/n,1/n,\dots,n/n\}.$$
\end{prob}
\end{frame}

\begin{frame}{Challenge!}
Suppose that $f(x) = x$, $\alpha(x) = x$, and $[a,b] = [0,1]$.\\
\begin{prob}
Prove that $\int_0^1 f d\alpha = 1/2$.
\end{prob}
\end{frame}

\begin{frame}{Challenge!}
Suppose that $f(x) = x^2$, $\alpha(x) = x$, and $[a,b] = [0,1]$.\\
\begin{prob}
Prove that $\int_0^1 f d\alpha = 1/3$.
Hint:
$$1^2 + 2^2 + 3^2 + \dots + n^2 = \frac{n(n+1)(2n+1)}{6}$$
\end{prob}
\end{frame}

\begin{frame}{Continuous implies integrable}
\begin{thm}
Suppose that $f$ is a continuous real-valued function on $[a,b]$.\\
Then $f$ is Riemann integrable on $[a,b]$.
\end{thm}
\end{frame}

\begin{frame}{Mean-Value Theorem for Integrals}
\begin{thm}
Suppose that $f$ is integrable with respect to $\alpha$ on $[a,b]$.\\
Let
$$m = \inf\{f(x): x\in [a,b]\},\quad\text{and}\quad M = \sup\{f(x): x\in [a,b]\}.$$
Then there exists $m\leq c\leq M$ with
$$\int_a^b f d\alpha = c(\alpha(b)-\alpha(a)).$$
In the special case $f$ is continuous, then $c = f(x_0)$ for some $x_0\in [a,b]$.
\end{thm}
\end{frame}

\begin{frame}{Fundamental Theorem of Calculus I}
Assume that $f$ is integrable with respect to $\alpha$ on $[a,b]$ and set\\
\pause
$$F(x) = \int_a^x fd\alpha,\quad x\in [a,b].$$
\begin{thm}[FTC I]
Suppose that $\alpha$ is monotone increasing on $[a,b]$.\\
Then for any $x\in (a,b)$ where $f$ is continuous and $\alpha$ is differentiable
$$F'(x) = f(x)\alpha'(x).$$
\end{thm}
\end{frame}

\begin{frame}{Fundamental Theorem of Calculus II}
\begin{thm}[FTC II]
Suppose that $f$ is Riemann integrable on $[a,b]$.\\
Also suppose $g$ is continuous on $[a,b]$ with $g'(x) = f(x)$ for all $x\in (a,b)$.\\
Then
$$\int_a^b f(x)dx = g(b)-g(a).$$
\end{thm}
\end{frame}


\end{document}


